\documentclass[10pt,a4paper]{amsart}
\usepackage[margin=19mm]{geometry}
\usepackage{amsmath,amssymb,mathtools}
\usepackage[scr=rsfs,cal=euler]{mathalfa}
\usepackage{xfrac}
\usepackage[unicode,hidelinks,bookmarks=false]{hyperref}
\newcommand{\ie}{i.e.\ }
\newcommand{\cf}{cf.\ }
\newcommand{\resp}{respectively\ }
\newcommand{\Subsection}{\subsection}
\providecommand{\nobreakdash}{\nobreak\hbox{-}}
\setlength{\emergencystretch}{2em}
\multlinegap0pt
\usepackage{amssymb}
\usepackage{enumerate}
\usepackage{tikz}
\newtheorem{lemma}{Lemma}
\title{Uncertainty propagation through trained multi-layer perceptrons: Exact analytical results}
\author{Andrew Thompson, Miles McCrory}
\date{Source: arXiv:2601.16830v1 (CC BY 4.0)}
\begin{document}
\maketitle

\noindent\emph{Source and licence.} Uncertainty propagation through trained multi-layer perceptrons: Exact analytical results. Version arXiv:2601.16830v1 (CC BY 4.0), distributed by arXiv under the Creative Commons Attribution 4.0 International licence, http://creativecommons.org/licenses/by/4.0/, as recorded in the arXiv licence metadata for this exact version and shown on the abstract page https://arxiv.org/abs/2601.16830v1. Full licence notice: \texttt{licenses/arxiv-2601.16830-CC-BY-4.0.txt}.\par\smallskip

\noindent\textit{Editorial scope.} Counted unit: the article's lemma truncated\_prob (source0.tex lines 253--258). The article's complete original proof of it is delivered with it (source0.tex lines 260--268, one \texttt{proof} environment of 9 lines). No mathematics is written, repaired or re-derived: the statement and the proof are the article's own lines, in the article's own notation, and no retained line is substituted or re-flowed.  Retained uncounted as the source-local context the proof needs: lines 132--132.  Nothing was omitted from any retained span, so the package carries no omission record.  No retained line contains a citation, so the wrapper carries no bibliography and no bibliography entry.  No retained line contains a figure environment, an image-inclusion command or drawing code, so no image is delivered and the package declares no asset dependency.  Editorial formatting only, in addition: an AMS \texttt{amsart} preamble that restates the article's own declarations and macros used by the retained lines, namely the environments \texttt{lemma} on separate counters, together with the standard packages those lines need.  The retained environments print in this document as Lemma 1 in order of appearance, whereas the article prints the same items with the numbering of its own full document.  The complete native source of the article as distributed in the arXiv e-print for this exact version is bundled unchanged as \texttt{sources/193264/source0.tex}.
\section{Analytical uncertainty propagation results}\label{analytical_results}

\begin{lemma}
Let $X$ be defined as in Section~\ref{analytical_results}. Then 
\begin{equation}\label{truncated_prob}\mathbb{P}(X_i > 0 \cap X_j > 0) = \Phi_2 \left( \frac{\mu_i}{\sigma_i}, \frac{\mu_j}{\sigma_j}; \rho_{ij} \right).
\end{equation}
\label{lem:conprob}
\end{lemma}

\begin{proof}
   
We have $\mathbb{P}(X_i > 0 \cap X_j > 0) = \mathbb{P}(W_i > 0 \cap W_j > 0)$ where $W \sim \mathcal{N}(\mu, \Sigma)$. Defining $V_i = \frac{W_i-\mu_i}{\sigma_i}, V_j = \frac{W_j-\mu_j}{\sigma_j}$, we have
\begin{equation}\label{eqn:standard_normal}\left[ {V_i \atop V_j} \right] \sim \mathcal{N} \left( \left[ {0 \atop 0} \right], \left[ {1 \atop \rho_{ij}} {\rho_{ij} \atop 1} \right] \right).\end{equation}
We then have
\begin{eqnarray}\mathbb{P}(W_i > 0 \cap W_j > 0) &=& \mathbb{P}(\sigma_i V_i + \mu_i > 0 \cap \sigma_j V_j + \mu_j > 0)\nonumber\\&=& \mathbb{P} \left( V_i > - \frac{\mu_i}{\sigma_i} \cap V_j > - \frac{\mu_j}{\sigma_j} \right)\nonumber\\&=&\mathbb{P} \left( V_i < \frac{\mu_i}{\sigma_i} \cap V_j < \frac{\mu_j}{\sigma_j} \right),\label{eqn:lemma1}\end{eqnarray}
where the last line follows by the symmetry of $V_i$ and $V_j$ about $0$. The result now follows by combining (\ref{eqn:standard_normal}) and (\ref{eqn:lemma1}).
    
\end{proof}

\end{document}
